\documentclass[11pt]{article}

\usepackage[margin=1in]{geometry}
\usepackage{booktabs}
\usepackage{tabularx}
\usepackage{array}
\usepackage{threeparttable}
\usepackage{caption}
\usepackage{xcolor}
\usepackage{hyperref}

\captionsetup{font=small, labelfont=bf}
\renewcommand{\arraystretch}{1.25}

\title{Validation Tooling in LTLsplitter: \\ Model Checkers and Realizability Tools}
\author{}
\date{\today}

\begin{document}
\maketitle

\section{Overview}

LTLsplitter validates an LLM-proposed decomposition of a FRET LTL requirement, and the
realizability of the resulting specification, using five external formal tools rather than
one. Three tools (\textsc{nuXmv}, \textsc{Spot}, \textsc{SPIN}) cross-check \emph{logical
equivalence} between the original requirement and the conjunction of its sub-formulas; two
tools (\textsc{Strix}, \textsc{Kind2}) independently check \emph{realizability} of the
specification against an adversarial environment. Using more than one tool per question is a
deliberate redundancy: each tool is built on a different algorithmic foundation and has
different blind spots, so agreement across tools is a stronger correctness signal than any
single tool's verdict. Table~\ref{tab:validation-tools} summarizes each tool; the differences
in approach are discussed in Section~\ref{sec:discussion}.

\begin{table}[htbp]
\centering
\footnotesize
\begin{threeparttable}
\caption{Model checkers and realizability tools used for specification validation in LTLsplitter.}
\label{tab:validation-tools}
\begin{tabularx}{\textwidth}{@{} l l >{\raggedright\arraybackslash}X >{\raggedright\arraybackslash}X >{\raggedright\arraybackslash}X @{}}
\toprule
\textbf{Tool} & \textbf{Underlying technique} & \textbf{Role in LTLsplitter} & \textbf{Benefits} & \textbf{Shortcomings} \\
\midrule
\textbf{nuXmv}
  & Symbolic model checking (BDD / SAT-based bounded model checking)
  & Requirement-set \emph{consistency} (ReqAn: can all requirements hold simultaneously?) and \emph{equivalence} checking of past-time LTL formulas
  & Natively understands past-time LTL (H/O/S/Y) -- FRET's native output form; typed variables (bounded int, enum) modeled exactly, not abstracted; one tool covers two of the four checks
  & Symbolic (BDD/SAT) engines can scale poorly on large or highly disjunctive variable domains; Linux binary distribution is fragile (missing shared-library versions on modern distros) \\
\textbf{Spot} (\texttt{ltlfilt})
  & Automata-theoretic equivalence (Büchi automaton language containment, via \texttt{--equivalent-to})
  & Independent \emph{equivalence}/tautology check of the future-time (ft) form of the same formulas
  & Fast, purpose-built equivalence test; accepts a quoted string as an opaque atomic proposition, so typed comparisons (e.g.\ \texttt{"x = 1"}) need no renaming scheme
  & Future-time LTL only -- cannot parse FRET's native past-time form, so it checks a translated form rather than the original requirement directly \\
\textbf{SPIN}
  & Explicit-state model checking (never-claim product automaton + exhaustive search)
  & Third, independently-implemented \emph{equivalence} check of the ft form, deliberately redundant with Spot
  & A genuinely different algorithm family (explicit-state vs.\ symbolic/automata) catches bugs the other two share by construction; produces concrete counterexample traces
  & Requires a hand-built, free-running Promela model per formula; the search can get stuck at depth~0 depending on the initial state chosen; needs a C compiler at runtime \\
\textbf{Strix}
  & Reactive synthesis (parity-game solving over $\omega$-automata)
  & GR(1)-style \emph{realizability}: does a system strategy exist that satisfies the spec against every adversarial input sequence?
  & Directly answers the two-player-game realizability question; mature, well-tested LTL synthesis engine
  & Parser accepts only bare-identifier boolean propositions -- typed/arithmetic atoms must be abstracted into fresh booleans, which can produce spurious \textsc{unrealizable} verdicts by dropping real correlations between values \\
\textbf{Kind2}
  & SMT-based $k$-induction over Lustre contracts (\texttt{CONTRACTCK})
  & Independent \emph{realizability} cross-check against Strix, via a contract-realizability analysis
  & Models int/real/enum variables natively (Lustre has real types) -- no abstraction gap, so it is not fooled by the artifact that can make Strix say \textsc{unrealizable}; separates environment assumptions from system guarantees
  & Has no native past-time operator, so past-time formulas must be hand-compiled into Lustre ``ghost'' variables recurring on \texttt{pre}; SMT/$k$-induction can fail to converge on complex contracts \\
\bottomrule
\end{tabularx}
\begin{tablenotes}
\footnotesize
\item GR(1) = Generalized Reactivity(1), a tractable fragment of LTL realizability (polynomial-time synthesis, vs.\ 2EXPTIME for general LTL).
\end{tablenotes}
\end{threeparttable}
\end{table}

\section{Discussion: what the table does not say explicitly}
\label{sec:discussion}

\textbf{Three different questions, not one.} ``Consistency'' (nuXmv/ReqAn) asks whether a set
of requirements can hold \emph{simultaneously} on some trace -- a satisfiability question.
``Equivalence'' (nuXmv, Spot, SPIN) asks whether two formulas denote the \emph{same} set of
traces -- used here to check that a decomposition is faithful to the original requirement, and
that no sub-formula is itself already the whole requirement in disguise (non-degeneracy).
``Realizability'' (Strix, Kind2) is a strictly harder, two-player question: does a
\emph{system strategy} exist that wins against every move an adversarial environment could
make -- this is undecidable in general and 2EXPTIME-complete for full LTL, which is why both
realizability tools restrict to a tractable fragment (GR(1) for Strix; contract-shaped
specifications for Kind2) rather than attacking arbitrary LTL directly.

\textbf{Redundancy is cross-validation, not duplication.} The three equivalence checkers agree
on the question but disagree on algorithm (symbolic BDD/SAT, automata-theoretic, explicit-state)
and on which temporal direction they accept (past-time vs.\ future-time LTL). A verdict that
only one tool reaches is treated with suspicion; agreement across algorithmically unrelated
tools is the actual correctness signal. The same logic applies to Strix vs.\ Kind2 for
realizability: they disagree on algorithm (game synthesis vs.\ SMT/$k$-induction) and, crucially,
on how precisely they model typed variables -- Kind2's native typing makes it the arbiter when
Strix's abstraction-driven \textsc{unrealizable} verdict is suspected to be an artifact rather
than a genuine synthesis failure.

\section{Supporting pipeline tools: FRET and Ogma/Copilot}

FRET and Ogma are part of the same end-to-end pipeline but are \emph{not} validation tools in
the sense of Table~\ref{tab:validation-tools} -- neither one checks a property of a
specification the way nuXmv, Spot, SPIN, Strix, or Kind2 do. Table~\ref{tab:pipeline-tools}
places them relative to the validation stage.

\begin{table}[htbp]
\centering
\footnotesize
\begin{tabularx}{\textwidth}{@{} l >{\raggedright\arraybackslash}X >{\raggedright\arraybackslash}X @{}}
\toprule
\textbf{Tool} & \textbf{What it is} & \textbf{Why it is not a validation tool} \\
\midrule
\textbf{FRET} (NASA)
  & A GUI for eliciting requirements in structured English (``FRETish'') and formalizing them into past-time LTL. It runs its own internal consistency/realizability checks (via NuSMV/Kind2) before export.
  & FRET's checks happen \emph{upstream} of LTLsplitter, on the original requirement, before import -- it is a requirements-elicitation and formalization tool, not something LTLsplitter invokes to validate a decomposition or synthesize a strategy. \\
\textbf{Ogma / Copilot} (NASA)
  & Ogma is a Haskell code-generation CLI that wraps Copilot, a stream-processing DSL compiling to hard-real-time C99. Given an already-validated LTL/Lustre/SMV specification, \texttt{ogma ros} generates a self-contained ROS2 package: a monitor node that subscribes to input topics, re-evaluates the properties, and publishes violations.
  & It performs \emph{runtime monitoring}, not design-time formal verification: it does not prove anything about the specification itself, it compiles an already-checked specification into executable code that watches a live system trace. It runs \emph{after} every check in Table~\ref{tab:validation-tools} has passed, not instead of them. \\
\bottomrule
\end{tabularx}
\caption{Supporting pipeline tools: requirements elicitation (FRET) and runtime monitor synthesis (Ogma/Copilot), shown against the design-time validation tools of Table~\ref{tab:validation-tools}.}
\label{tab:pipeline-tools}
\end{table}

In short: FRET sits \emph{before} validation (it produces the LTL that gets split and checked),
the five tools in Table~\ref{tab:validation-tools} sit \emph{during} validation (they check the
split is faithful and the result is realizable), and Ogma/Copilot sits \emph{after} validation
(it turns the now-trusted specification into a deployed runtime monitor). None of the three
stages substitutes for another: FRET's own NuSMV/Kind2 checks do not know about the LLM-proposed
split, and Ogma does not check correctness at all -- it assumes the specification handed to it
is already correct.

\end{document}
